\documentclass[../../../include/open-logic-section]{subfiles}

\begin{document}

\olfileid{sth}{replacement}{limofsize}
\olsection{محدودیتِ اندازه}

شاید رایج‌ترین کوشش برای ارائهٔ توجیهی «درونی» از جایگزینی بر مفهوم زیر
تکیه کند:
\begin{enumerate}
	\item[] \limofsize. هر چند چیز یک مجموعه تشکیل می‌دهند، مشروط بر
	اینکه تعدادشان بیش از حد زیاد نباشد.
\end{enumerate}
این اصل بی‌درنگ جایگزینی را موجه می‌سازد. زیرا مجموعه‌ای که با جایگزینی
تشکیل می‌شود نمی‌تواند از مجموعهٔ مبدأ خود پرشمارتر باشد. به‌بیان دقیق،
فرض کنید با جایگزینی مجموعهٔ
$\funimage{\tau}{A}=\Setabs{\tau(x)}{x\in A}$ را تشکیل می‌دهید. نگاشتی
که هر~$x\in A$ را به~$\tau(x)$ می‌فرستد، از~$A$ بر
$\funimage{\tau}{A}$ پوشاست؛ بنابراین تصویر، در معنای مقایسه با نگاشت
پوشا، از~$A$ پرشمارتر نیست. پس اگر~!!{element}های~$A$ آن‌قدر پرشمار
نبودند که نتوانند مجموعه‌ای تشکیل دهند، تصاویرشان نیز آن‌قدر پرشمار
نیستند که نتوانند~$\funimage{\tau}{A}$ را تشکیل دهند.

دشواری آشکارِ توسل به~\limofsize{} برای توجیه جایگزینی آن است که تاکنون
هیچ اصلی مانند~\limofsize{} وضع نکرده‌ایم. افزون بر این، هنگامی که در
\crefrange{sth:story::chap}{sth:z::chap} روایت خود را دربارهٔ تصور
انباشتی-تکراریِ مجموعه بیان کردیم، هیچ چیز حتی \emph{اشاره‌ای} به
\limofsize{} نداشت. درست به همین سبب بود که بولوس زمانی نوشت: «شاید بتوان
نتیجه گرفت که دست‌کم دو اندیشه در ``پسِ'' نظریهٔ مجموعه‌ها وجود دارد»
\citeyearpar[p.~19]{Boolos1989}. از یک سو، قرار است اندیشه‌های پیرامون
تصور انباشتی-تکراریِ مجموعه، $\Z$ را موجه سازند. از سوی دیگر، قرار است
\limofsize{} جایگزینی را موجه سازد.

اما مسئله فقط این نیست که تاکنون دربارهٔ~\limofsize{} \emph{سکوت}
کرده‌ایم. بلکه مسئله آن است که~\limofsize{} (با صورت‌بندی کنونی) چندان
با تصور انباشتی-تکراریِ مجموعه سازگار به نظر نمی‌رسد. زیرا هیچ اشاره‌ای
به اندیشهٔ تشکیل مجموعه‌ها در \emph{مراحل} ندارد.

با توجه به تاریخ این «دو اندیشه» (یعنی تصور انباشتی-تکراریِ مجموعه و
\limofsize)، این امر چندان شگفت‌آور نیست. این «دو اندیشه» در نهایت دو
برنامهٔ کاملاً متفاوت برای جلوگیری از پارادوکس‌های نظریهٔ مجموعه‌ها
هستند. تصور انباشتی-تکراریِ مجموعه، به‌بیان تقریبی، پارادوکس راسل را
چنین مسدود می‌کند: \emph{هرگز نباید انتظار می‌داشتیم مجموعهٔ راسل وجود
داشته باشد، زیرا در هیچ مرحله‌ای «تشکیل» نمی‌شد}. در مقابل، قرار است
\limofsize{} مجموعهٔ راسل را چنین کنار بگذارد: \emph{هرگز نباید انتظار
می‌داشتیم مجموعهٔ راسل وجود داشته باشد، زیرا بیش از حد بزرگ می‌بود}.

اگر مطلب را چنین بیان کنیم، باید صریح باشیم: \limofsize{} به‌عنوان پاسخی
به پارادوکس‌ها به توجیهی \emph{بسیار} بیشتر نیاز دارد. برای نمونه، این
صورت از پارادوکس راسل را در نظر بگیرید: \emph{هیچ سگ پاگ دقیقاً همان
پاگ‌هایی را بو نمی‌کشد که خودشان را بو نمی‌کشند}
(\olref[sth][story][rus]{sec} را ببینید). اگر بپرسید «چرا چنین پاگی وجود
ندارد؟»، پاسخ خوبی نیست که بگویند چنین پاگی ناچار بود تعداد بیش از حدی
پاگ را بو بکشد. پس چرا اینکه مجموعهٔ راسل برای وجودداشتن «بیش از حد
بزرگ» می‌بود، باید توضیح شهودی خوبی برای نبودِ آن باشد؟

خلاصه آنکه، اگر انگیزهٔ «شهودیِ»~\limofsize{} اندکی سردرگمتان می‌کند،
جای سرزنش نیست.

\end{document}
